% ============================================================
% paper/sections/09_parent_coupling_k7.tex
% Parent Coupling: Constraints from the K7 Continuum (finalized)
% + FWHQ-01 foundational hardening additions
% ============================================================

\section{Parent Coupling: Constraints from the K7 Continuum}

\subsection{Non-Isolated Nature of the Enterprise Continuum}

The enterprise continuum $K_{\mathrm{EA}}$ does not exist as an isolated
formal object.
It is embedded within a parent continuum, denoted $K_7$, representing
the broader socio-technical context in which enterprise states must be
admissible.

Coupling to the parent continuum does not prescribe enterprise behavior,
goals, or evolution.
It constrains admissibility by delimiting the external conditions under
which $K_{\mathrm{EA}}$ may exist, persist, and remain diagnosable.

Parent coupling therefore functions as a boundary condition, not as a
causal or explanatory model of enterprise behavior.

\subsection{External Fields as Exogenous Constraints}

According to \texttt{ETS.K\_EA.v1.1}, a finite set of external fields is
recognized as admissible couplings from $K_7$.
These fields act as exogenous constraints rather than control variables
or drivers of internal dynamics:

\begin{itemize}
  \item \textbf{Law}: conditions of legal existence, liability, and
        enforceability.
  \item \textbf{Regulation}: sector-specific and cross-sector constraint
        regimes.
  \item \textbf{Market}: demand conditions, competitive pressure, and
        economic feasibility.
  \item \textbf{Technology}: available technical substrate and material
        limits.
  \item \textbf{Norms}: socially and organizationally enforced
        expectations.
\end{itemize}

These fields restrict the admissible region of
$\Omega_{\mathrm{EA}}$ without directly specifying enterprise-internal
states, trajectories, or targets.

\subsection{Indirect Effects on Enterprise Primitives}

External fields affect the enterprise continuum only through their
indirect interaction with existing primitives.
Direct control of enterprise-internal variables is not permitted within
the model.

In particular:
\begin{itemize}
  \item legal and regulatory fields constrain existence and termination
        thresholds,
  \item market conditions constrain resource-related potentials and
        inertia-related thresholds,
  \item technological fields constrain admissible structural and
        informational axes,
  \item normative fields constrain observability and semantic alignment.
\end{itemize}

All effects are mediated through $\Theta_{\mathrm{EA}}$,
$P_{\mathrm{EA}}$, $A_{\mathrm{EA}}$, or observability conditions.
No external field acts as an independent primitive, operator, or flow
within $K_{\mathrm{EA}}$.

\subsection{Compatibility Constraints}

The Executable Theory Specification enforces explicit compatibility
constraints between $K_{\mathrm{EA}}$ and its parent continuum:

\begin{itemize}
  \item $\Omega_{\mathrm{EA}} \subseteq \Omega_{7}$: admissible enterprise
        states must remain within the admissible state space of the
        parent continuum.
  \item No direct external axis control: parent fields may not set,
        override, or optimize values of $A_{\mathrm{EA}}$.
  \item No direct $K_8 \rightarrow K_{\mathrm{EA}}$ coupling: higher-level
        continua may affect $K_{\mathrm{EA}}$ only through $K_7$.
\end{itemize}

Violation of any compatibility constraint triggers termination of
diagnosis.
Such termination reflects loss of definability, not empirical collapse
or failure.

\subsection{Interpretation Limits}

Parent coupling specifies how external conditions delimit the domain of
admissible enterprise diagnosis.
It does not explain the origin of those conditions, predict their
evolution, quantify causal influence, or prescribe responses.

No sociological, economic, or institutional theory is implied by this
coupling structure.
At all levels, the model remains diagnostic, non-prescriptive, and
executable-closed.

%%%End of Document%%%
